% ============================================================
%  Преамбула: пакеты, шрифты, геометрия, оформление
%  Стандарт оформления КВР МФТИ ФПМИ:
%    - Times New Roman 14 pt, интервал 1.5
%    - Поля: левое 30 мм, правое 15 мм, верх/низ 20 мм
%    - Абзацный отступ 1.25 см
%    - Нумерация страниц снизу по центру
%    - Заголовки глав с новой страницы
% ============================================================

% --- Языковая поддержка (XeLaTeX/LuaLaTeX) ---
\usepackage{polyglossia}
\setdefaultlanguage{russian}
\setotherlanguage{english}

% --- Шрифты ---
\usepackage{fontspec}
\defaultfontfeatures{Ligatures=TeX}
\setmainfont{Times New Roman}
\setsansfont{Arial}
\setmonofont{Courier New}
\newfontfamily\cyrillicfont{Times New Roman}
\newfontfamily\cyrillicfontsf{Arial}
\newfontfamily\cyrillicfonttt{Courier New}

% --- Геометрия ---
\usepackage[a4paper,
            left=30mm, right=15mm,
            top=20mm, bottom=20mm,
            footskip=10mm]{geometry}

% --- Интервал 1.5 ---
\usepackage{setspace}
\onehalfspacing

% --- Абзацный отступ ---
\usepackage{indentfirst}
\setlength{\parindent}{1.25cm}
\setlength{\parskip}{0pt}

% --- Математика ---
\usepackage{amsmath, amssymb, amsthm, mathtools}
\usepackage{bm}

% --- Графика ---
\usepackage{graphicx}
\usepackage{tikz}
\usetikzlibrary{positioning, arrows.meta, calc, shapes.geometric}
\usepackage{float}
\graphicspath{{figures/}}

% --- Таблицы ---
\usepackage{array}
\usepackage{booktabs}
\usepackage{tabularx}
\usepackage{multirow}
\usepackage{longtable}

% --- Подписи к рисункам и таблицам ---
\usepackage[labelsep=endash, font=small, justification=centering]{caption}
\usepackage{subcaption}
\captionsetup[figure]{name=Рисунок}
\captionsetup[table]{name=Таблица}

% --- Гиперссылки ---
\usepackage[hidelinks, unicode=true]{hyperref}
\usepackage{cleveref}
% Явно задаём имена для всех типов ссылок, чтобы cleveref не пытался
% синтезировать «Глава» через babel-макросы \cyrg\cyrl\cyra\cyrv,
% которых нет под polyglossia.
\crefname{chapter}{гл.}{гл.}
\Crefname{chapter}{Глава}{Главы}
\crefname{section}{разд.}{разд.}
\Crefname{section}{Раздел}{Разделы}
\crefname{subsection}{разд.}{разд.}
\Crefname{subsection}{Раздел}{Разделы}
\crefname{figure}{рис.}{рис.}
\Crefname{figure}{Рисунок}{Рисунки}
\crefname{table}{табл.}{табл.}
\Crefname{table}{Таблица}{Таблицы}
\crefname{equation}{}{}
\Crefname{equation}{Уравнение}{Уравнения}
\crefname{proposition}{утв.}{утв.}
\Crefname{proposition}{Утверждение}{Утверждения}

% --- Списки ---
\usepackage{enumitem}
\setlist{nosep, leftmargin=*}

% --- Заголовки глав/секций (ГОСТ-стиль ФПМИ) ---
\usepackage{titlesec}
\titleformat{\chapter}[block]
  {\normalfont\Large\bfseries\centering}
  {\thechapter}{1ex}{}
\titlespacing*{\chapter}{0pt}{0pt}{2ex}

\titleformat{\section}[block]
  {\normalfont\large\bfseries}
  {\thesection}{1ex}{}
\titleformat{\subsection}[block]
  {\normalfont\normalsize\bfseries}
  {\thesubsection}{1ex}{}

% --- Нумерация страниц ---
\usepackage{fancyhdr}
\pagestyle{plain}

% --- Библиография через biblatex (ГОСТ) ---
\usepackage[backend=biber,
            style=gost-numeric,
            sorting=none,
            language=auto,
            autolang=other,
            doi=true,
            isbn=false,
            url=true]{biblatex}

% --- Теоремы, определения ---
\newtheorem{theorem}{Теорема}
\newtheorem{lemma}{Лемма}
\newtheorem{proposition}{Утверждение}
\theoremstyle{definition}
\newtheorem{definition}{Определение}
\newtheorem{example}{Пример}
\theoremstyle{remark}
\newtheorem{remark}{Замечание}

% --- Удобные сокращения ---
\newcommand{\R}{\mathbb{R}}
\newcommand{\E}{\mathbb{E}}
\newcommand{\Q}{\mathbb{Q}}
\newcommand{\Prob}{\mathbb{P}}
\newcommand{\Var}{\operatorname{Var}}
\newcommand{\Cov}{\operatorname{Cov}}
\newcommand{\Corr}{\operatorname{Corr}}
\newcommand{\dd}{\,\mathrm{d}}
\DeclareMathOperator{\sech}{sech}

% --- Заглушка-плейсхолдер для графиков (зелёный квадрат) ---
% Использование: \placeholderfig{ширина}{высота}{подпись}{метка}
\newcommand{\placeholderfig}[4]{%
  \begin{figure}[H]
    \centering
    \begin{tikzpicture}
      \fill[green!55!black!40] (0,0) rectangle (#1,#2);
      \draw[green!40!black, thick] (0,0) rectangle (#1,#2);
      \node[font=\small\itshape, text=black!70]
            at (#1/2, #2/2) {Место для графика};
    \end{tikzpicture}
    \caption{#3}
    \label{#4}
  \end{figure}}

% Заглушка для таблицы
\newcommand{\placeholdertab}[2]{%
  \begin{table}[H]
    \centering
    \caption{#1}\label{#2}
    \begin{tabular}{|c|c|c|c|}
      \hline
      \multicolumn{4}{|c|}{\textit{Место для таблицы --- будет заполнено}}\\
      \hline
      --- & --- & --- & --- \\
      \hline
      --- & --- & --- & --- \\
      \hline
    \end{tabular}
  \end{table}}

% Заглушка-абзац для будущего текста
\newcommand{\todoblock}[1]{%
  \par\medskip
  \noindent\fbox{\parbox{0.96\linewidth}{\textit{TODO:} #1}}
  \par\medskip}
